% demo.tex — HUST Beamer Theme 演示文件
% 编译方式: xelatex demo.tex (需编译两次以获得正确页码和导航)

\documentclass[aspectratio=169]{beamer}
% 4:3 比例请用: \documentclass{beamer}

\usepackage{ctex}  % 中文支持
\usetheme{HUST}

% === 设置 Logo 路径 ===
\hustlogo{xiaohui.png}
\physicslogo{物理学院 logo.png}
\gravitylogo{国家精密重力测量科学中心LOGO-左右带英文.png}

% === 演讲信息 ===
\title{HUST Beamer Theme Demo}
\subtitle{华中科技大学 Beamer 主题演示}
\author{Your Name}
\institute{华中科技大学}
\date{\today}

\begin{document}

% ===== 首页（必须使用 plain 选项以隐藏侧边栏）=====
\begin{frame}[plain]
  \titlepage
\end{frame}

% ===== 目录 =====
\begin{frame}{目录}
  \tableofcontents
\end{frame}

% ===== 第一节 =====
\section{研究背景}

\begin{frame}{脉冲星计时阵列简介}
  \begin{itemize}
    \item 脉冲星计时阵列（PTA）是探测纳赫兹引力波的重要手段
    \item 红噪声是 PTA 数据分析中的主要挑战之一
    \item 传统方法在去除红噪声时存在局限性
  \end{itemize}

  \begin{block}{主要困难}
    红噪声与引力波信号在频域上高度重叠，分离难度大。
  \end{block}
\end{frame}

% ===== 第二节 =====
\section{方法}

\begin{frame}{新方法概述}
  \begin{itemize}
    \item 提出基于贝叶斯框架的红噪声建模方案
    \item 引入自适应频率分辨方法
    \item 结合 MCMC 采样实现高效参数估计
  \end{itemize}

  \begin{alertblock}{注意}
    该方法需要足够长的观测时间基线（$>$ 10 年）。
  \end{alertblock}
\end{frame}

% ===== 第三节 =====
\section{结果}

\begin{frame}{主要结果}
  \begin{enumerate}
    \item 噪声去除效果显著提升（信噪比提高 30\%）
    \item 计算效率相比传统方法提高约 5 倍
    \item 在模拟和真实 IPTA DR2 数据上均验证有效
  \end{enumerate}

  \begin{exampleblock}{示例}
    在 NANOGrav 15yr 数据集上的测试结果表明，
    新方法能够更准确地恢复注入的引力波信号。
  \end{exampleblock}
\end{frame}

% ===== 第四节 =====
\section{总结与展望}

\begin{frame}{总结与展望}
  \begin{itemize}
    \item 本工作提出了一种新的红噪声去除方法
    \item 在模拟和真实数据上均取得了良好效果
    \item 未来计划：
      \begin{itemize}
        \item 进一步优化算法性能
        \item 拓展至空间引力波探测器数据分析
        \item 与国际 PTA 合作组开展联合分析
      \end{itemize}
  \end{itemize}

  \vfill
  \centering
  {\Large\color{HUSTBlue} 谢谢！}
\end{frame}

\end{document}
